\documentclass[11pt,a4paper]{article}
\usepackage[T1]{fontenc}
\usepackage[utf8]{inputenc}
\usepackage[margin=2.3cm]{geometry}
\usepackage{lmodern}
\usepackage{hyperref}
\usepackage{booktabs}
\usepackage{enumitem}
\setlist{noitemsep,topsep=4pt}
\setlength{\parindent}{0pt}
\setlength{\parskip}{6pt}

\title{Chess Club Tournament \& Rating Manager\\Developer Manual}
\author{Yahya El Idrissi\\THGA Bochum -- Introduction to Database Management Systems}
\date{September 2026}

\begin{document}
\maketitle
\tableofcontents
\newpage

\section{System Architecture}
\begin{verbatim}
tkinter desktop frontend (.deb)
        -> HTTP + X-API-Key
FastAPI container
        -> psycopg
PostgreSQL 16 container
\end{verbatim}
Docker Compose orchestrates both backend services. The frontend is not containerized and never accesses PostgreSQL directly. At startup it asks for the API URL and \texttt{X-API-Key}; these values are used for the current process.

\section{Database Model}
The schema in \texttt{db/01\_schema.sql} contains seven tables: \texttt{player}, \texttt{opening}, \texttt{tournament}, \texttt{registration}, \texttt{round}, \texttt{game}, and \texttt{rating\_history}.

Important integrity rules include primary and foreign keys, the composite registration key \texttt{(player\_id,tournament\_id)}, unique tournament round numbers, valid result values, distinct White/Black players, and controlled tournament status values. \texttt{db/02\_seed.sql} contains deterministic reference/sample data used by CI.

\section{Backend API and Business Logic}
The FastAPI entry point is \texttt{api/app/main.py}. \texttt{/health} is public; application routes require \texttt{X-API-Key}. The API supports players, tournaments, registrations, games, rating history, and standings.

Multi-row writes are transactional. Player creation writes the player and first history row. Tournament creation writes the tournament and all rounds. Game recording validates the round and registrations, inserts the game, updates both Elo values, and inserts two history rows; errors roll the transaction back.

Elo calculation is isolated in \texttt{api/app/elo.py} with K=20. Standings use SQL common table expressions. A win is 1 point, a draw 0.5, and a loss 0. Full Buchholz is the sum of opponents' tournament points.

\section{Frontend}
The tkinter frontend contains six tabs: Players, Standings, Rating History, Record Game, Register Player, and Create Tournament.

\texttt{api\_client.py} is the frontend HTTP layer. \texttt{main.py} first opens a startup connection dialog for API URL and key, validates the connection by calling the API, and then shows the main window. Successful writes refresh dependent read views.

\section{Build, Run and Package}
Create configuration and start the backend:
\begin{verbatim}
cp .env.example .env
docker compose up -d --build
curl http://127.0.0.1:8000/health
\end{verbatim}
Start the source frontend:
\begin{verbatim}
PYTHONPATH="$PWD/frontend" python3 -m chess_club_fe.main
\end{verbatim}
For Debian packaging:
\begin{verbatim}
sh frontend/debian/build-deb.sh
\end{verbatim}
The generated package is \texttt{frontend/debian/chess-club-manager\_1.0.0\_all.deb}. The installed frontend uses the same startup connection dialog; the backend remains in Docker Compose.

\section{Tests, CI and Documentation Build}
For local API tests:
\begin{verbatim}
python3 -m venv api/.venv
api/.venv/bin/pip install -r api/requirements.txt
api/.venv/bin/pip install -r api/requirements-dev.txt
cd api
.venv/bin/pytest tests -v
\end{verbatim}
Build all documentation PDFs with:
\begin{verbatim}
make all
\end{verbatim}
The Makefile writes the proposal, User Manual, and Developer Manual PDFs to \texttt{out/}.

GitHub Actions verifies the database schema and seed, API tests, Docker Compose startup of PostgreSQL + FastAPI, Debian package creation, and Makefile-driven LaTeX builds. On a pushed version tag matching \texttt{v*}, CI creates a GitHub Release and attaches the generated PDFs.

\section{How to Add a New Function}
A new feature should follow the existing layers:
\begin{enumerate}
  \item Decide whether persistent data or schema changes are required; preserve keys, constraints, and indexes.
  \item Add or extend a FastAPI route and reuse the API-key dependency for application data.
  \item Use a transaction whenever one action changes multiple rows.
  \item Extend \texttt{frontend/chess\_club\_fe/api\_client.py}; never access PostgreSQL from tkinter.
  \item Add the GUI control or tab and refresh affected views after writes.
  \item Add tests for success, authentication, and important invalid cases.
  \item Run pytest, Compose smoke testing, packaging, and documentation CI.
  \item Update the User Manual, Developer Manual, and README.
\end{enumerate}

A tournament-close feature, for example, would require a backend state transition, frontend API-client method, GUI action, tests for allowed/disallowed transitions, and documentation.

\section{Possible Future Extensions}
Reasonable future work includes tournament closing controls, edit/delete workflows, automatic pairing generation, richer opening/player analytics, improved round selection, and role-based users. These are intentionally outside the delivered scope.

\section{Reflection}
Separating the database, API, and GUI made the project easier to test and reason about: database constraints can be verified independently, API behavior can be tested before GUI work, and frontend code does not contain SQL. The most consistency-sensitive operation is game recording because one user action affects the game table, two current Elo values, and two history rows; the transaction prevents partial updates.

The deployment model also became part of the design. Containerizing PostgreSQL and FastAPI together makes the backend reproducible, while keeping the tkinter frontend outside Docker matches the required desktop installation model.

\section{Sources and Reuse}
Course patterns were reused and adapted where appropriate, especially the PostgreSQL/FastAPI separation, \texttt{X-API-Key} protection, Docker Compose backend structure, and tkinter frontend approach. These follow the DBMS lecture/exercise material around FastAPI, Docker Compose, API-key protection, and frontend packaging. Project-specific schema design, Elo logic, standings/Buchholz queries, tournament workflow, tests, and documentation were developed for this chess-club application.

Any future third-party code should be documented with its exact source and the parts that were adapted.

\end{document}
